% SPDX-License-Identifier: CC-BY-SA-4.0
% Author: Matthieu Perrin
% Part: 
% Section: 
% Sub-section: 
% Frame: 

\begingroup

\begin{frame}{Modèle de communication}

  \begin{block}{Modèle à mémoire partagée}
    \begin{itemize}
    \item Communication par des \structure{instructions} sur des \alert{registres partagés} 
      \begin{itemize}
      \item \lstinline{get}, \lstinline{set},
      \item \lstinline{compareAndSet}, \lstinline{getAndSet}, \lstinline{getAndAdd}...
      \end{itemize}
    \item On suppose une mémoire \structure{infinie} organisée en \alert{tas}
      \begin{itemize}
      \item allocation dynamique par un mot-clé \lstinline{new},
      \item présence d'un \structure{garbage collector}
      \end{itemize}
    \item L'ensemble des registres est noté \alert{$\textsc{Reg}$}
    \item L'ensemble des valeurs qu'ils peuvent prendre est noté \alert{$\textsc{Val}$}
    \end{itemize}
  \end{block}

  \pause
  
  \begin{alertblock}{Hypothèse simplificatrice}
    \begin{itemize}
    \item Toutes les instructions sont \alert{atomiques}
      \begin{itemize}
      \item un appel à une instruction atomique est un pas d'un processus
      \item en pratique, utiliser \lstinline{volatile} ou \lstinline{atomic*}
      \end{itemize}
    \end{itemize}
  \end{alertblock}

\end{frame}

\endgroup
\endinput

